\section{Reproducibility audit and protocol limits}
\label{sec:reproext}

\subsection{Source records and checksums}
The original Veltri split is distributed as six FASTA files in the authors'
\texttt{original-dataset} folder. The updated production dataset is a ZIP
archive linked from the AMP Scanner news page. The script
\texttt{studies/fetch\_published\_benchmarks.py} retrieves the files, checks
that each original training/test split has 712 positives and 712 decoys,
that the tune split has 354 each, and that the updated dataset has 2{,}021
positives and 2{,}021 decoys. No exact-sequence duplicates or cross-class exact overlaps were found in
the six original split files or within the updated positive and negative
sets; shared motifs and distant homology are not ruled out. In fact, a separate
exact-10-mer audit of the full Feb2020 stratified CV folds finds that
646/4{,}042 held-out sequences (15.98\%) share at least one exact decamer
with a training sequence. This is a limitation of the published random-split
style benchmark and of our replication, not a leakage-free generalization
estimate; see \texttt{results/study16\_homology\_audit.json}.
SHA-256 checksums for each downloaded FASTA
are in \texttt{results/benchmark\_sources.json}.

The original partition was used for studies 10 and 11. Studies 13 and 14
were run on the updated production dataset. Study 13 is deliberately kept
as a \emph{non-comparable ablation}: it used a 150-residue cap, retaining
2{,}013 positives and 2{,}020 decoys, and removed unknown amino acid symbols
from two decoys. Study 14 retains all 4{,}042 source records, represents X
as an unknown residue (zero one-hot/physicochemical row), and uses a
200-residue cap, above the observed maximum of 183. It alone supplies the
full-record comparison.

\subsection{Reproduction commands}
\begin{verbatim}
python studies/fetch_published_benchmarks.py
python studies/study10_veltri.py
python studies/study11_veltri_stack.py
python studies/study13_feb2020_cv.py # filtered ablation only
python studies/study14_feb2020_full.py # complete dataset
python studies/study12_candidates.py
python studies/make_sota_section.py
python studies/make_figs2.py
python -m pytest -q tests
cd paper && pdflatex -interaction=nonstopmode main.tex
\end{verbatim}

The complete 10-fold CNN/GNN/RF experiment takes about 10--20 minutes on
the stated two-CPU sandbox. It is not part of the hermetic test suite; tests
use tiny fixtures and never access the network. The random-forest fit is
deterministic from its seed. The convolutional training helper draws a
permutation once and reuses it across epochs; this is a design choice that
limits stochasticity and should be tested against reshuffling each epoch.

\subsection{Inference and statistical limits}
The 10-fold standard deviation measures variation across \emph{our} folds,
not uncertainty in the gap from the published AMP Scanner result: the authors'
per-fold scores and fold assignments are unavailable. The difference of two
published means cannot be assigned a paired $p$-value without those records.
A numerical lead of less than one percentage point in accuracy or a few
thousandths in AUC may disappear under resampling, preprocessing changes,
or alternative hyperparameters. Our bootstrap intervals in the UniProt
experiments resample individual sequences, not homology clusters, and may
understate uncertainty for correlated families. No independent potency
regression or clinical assay was run.

\subsection{Candidate ranking limits}
All candidate sequences in study 12 were selected after looking at the
study-08 model scores. The selection is conditional and the resulting scores
cannot be treated as unbiased validation metrics. P0ACW4 and P0ACW5 have
identical sequences and are collapsed before naming the unique candidates.
The sequence of P0ACW4 has an N-terminal signal peptide annotated at residues
1--20; a predictor trained on AMP labels may learn secretion or short-protein
features instead of antibacterial activity. The designed variant was
optimized on the same ensemble that scored it, so its higher score cannot
be taken as an independent confirmation. Pairwise 3-mer Jaccard against a
finite corpus is a useful triage feature, not proof of sequence novelty
in all public databases. Antimicrobial efficacy would require synthesized
purified peptide, growth inhibition across concentrations with positive and
negative controls, replication, and toxicity assessment.
